\documentclass{article}

\usepackage{PRIMEarxiv}

\usepackage[utf8]{inputenc} % allow utf-8 input
\usepackage[T1]{fontenc}    % use 8-bit T1 fonts
\usepackage{hyperref}       % hyperlinks
\usepackage{url}            % simple URL typesetting
\usepackage{booktabs}       % professional-quality tables
\usepackage{amsfonts}       % blackboard math symbols
\usepackage{amsmath}        % \text in equations
\usepackage{nicefrac}       % compact symbols for 1/2, etc.
\usepackage{microtype}      % microtypography
\usepackage{lipsum}
\usepackage{fancyhdr}       % header
\usepackage{graphicx}       % graphics
\graphicspath{{media/}}     % organize your images and other figures under media/ folder

%Header
\pagestyle{fancy}
\thispagestyle{empty}
\rhead{ \textit{ }} 

% Update your Headers here
\fancyhead[LO]{Jounit: a settlement layer for autonomous agents, machines and robots}
% \fancyhead[RE]{Firstauthor and Secondauthor} % Firstauthor et al. if more than 2 - must use \documentclass[twoside]{article}



  
%% Title
\title{Jounit: A Proof-of-Work Settlement Layer for Humans, Machines, Robots
and AI Agents\\
{\large Design, Security Analysis and Implementation Report of the JNT Protocol}
%%%% Cite as
%%%% Update your official citation here when published 
\thanks{\textit{\underline{Citation}}: 
\textbf{Gadila, S. Jounit: A Proof-of-Work Settlement Layer for Humans, Machines,
Robots and AI Agents. Preprint, 2026. DOI:000000/11111.}} 
}

\author{
  Swamy Gadila \\
  Jounit Protocol (community-owned, open source) \\
  \texttt{https://github.com/swamy18/Jounit} \\
}


\begin{document}
\maketitle


\begin{abstract}
Jounit is an open-source cryptocurrency and settlement protocol designed around a
single question that general-purpose blockchains do not answer well: how should a
machine be paid for work that a third party can verify without re-running it? This
paper describes the protocol end to end --- its proof-of-work layer, its task
marketplace, its wallet and agent-authorisation layer, and its security
engineering --- from the perspective of its sole author and implementer. The
consensus core is a linear UTXO chain whose fork choice is heaviest cumulative
work, whose proof of work is a three-stage hybrid (a SHA-256 chain, an
integer-only neural network seeded from the epoch identifier, and a binding
hash), and whose difficulty is adjusted by a multiplicative, windowed rule. Above
it sits a Layer-2 task economy in which submitters escrow payment, workers claim
and execute, anyone may challenge a result, and a stake-weighted committee of
BLS-signing validators adjudicates disputes; a second Layer-2 mechanism moves
value back to Layer 1 through an inclusion-proof-authorised, permissionlessly
crankable exit. Above both sits the layer this paper's later sections are
about: agent authority and credit as \emph{consensus} state rather than as
node-local bookkeeping --- a grant published on chain that binds its owner to a
named spender and a named agent, whose caps, allow-lists, window and nonce the
chain enforces, and a loan whose lender signs the terms, whose principal
really moves, and whose collateral answers for the debt. The paper documents the
trust boundaries explicitly (what is trusted, what is not, and why reputation is
deliberately excluded from consensus),
the concrete security mechanisms that were built and the audit-driven improvements
that followed them, the wallet and key-management design, the operational tooling
(HTTP API, keyless MCP surface, worker SDK, command line), the open-source and
governance
posture, and the limitations that remain before mainnet deployment.
\end{abstract}


% keywords can be removed
\keywords{Proof of Work \and Agent Payments \and Escrow \and Task Markets \and
Autonomous Agents \and Consensus Security \and Post-Quantum Signatures}


\section{Introduction}
\label{sec:intro}

\subsection{The problem}

Every general-purpose blockchain was designed for a world in which the party
holding a key is a person, a company, or a device operated by one. That
assumption breaks in two directions when software begins to act on a principal's
behalf. First, an autonomous agent that receives a payment must be able to prove
what it did and to be paid for it without a human approving every step. Second,
a principal that delegates spending authority to such an agent needs a way to
bound that authority --- by amount, by recipient, by time window --- and to
revoke it, without re-authorising each transaction by hand.

The second failure mode is the more consequential one, and it is not specific to
cryptocurrencies. Agent frameworks are proliferating rapidly \cite{sutton2019bitter,
yao2023survey,wang2023autogpt,yang2024sweagent}, and robotic platforms are
standardising on shared middleware such as the Robot Operating System
\cite{quan2021ros}. Both create economic actors --- software that consumes
resources on somebody's behalf --- while the payment rails remain
institution-centric. A language-model agent that rents compute, a warehouse
robot that consumes electricity, and a research pipeline that buys data are all
in the same position: they can produce a receipt, but they cannot enforce one.

Payment is only the visible part of the problem. The deeper issue is
\emph{verifiability}. If a machine claims it produced a result, the payer has
three unattractive options: trust the claim, re-run the work, or hire an
arbitrator they both accept. Re-running is often impossible (the work may have
consumed a one-shot resource), and a shared arbitrator is a trusted party --- the
exact thing a blockchain exists to avoid.

\subsection{The claim}

Jounit (ticker \texttt{JNT}) is an attempt to build the missing settlement
primitive directly into a cryptocurrency: a chain that can act as escrow and
arbiter for work performed by machines. The economic claim is deliberately
narrow and testable:

\begin{quote}
If a machine can produce a result that anyone else can check without re-running
it, then that machine can be paid for that result without either side trusting
the other.
\end{quote}

Everything in the protocol is arranged to make that sentence true in practice:
the escrow holds the payer's money before the work starts, the verifier is a
deterministic function of committed inputs, the adjudicator is a quorum of
stake-weighted validators rather than a named party, and the payout path is
ordinary on-chain value movement rather than a bespoke accounting entry.

\subsection{Contributions}

This paper is a design-and-implementation report rather than a theoretical
result. Its contributions are:

\begin{itemize}
\item \textbf{A concrete hybrid proof of work.} A three-stage puzzle that keeps
the security of a SHA-256 chain while attaching a bounded, epoch-seeded neural
stage, with a registry-enforced work-share band and integer-only arithmetic
throughout \cite{zawy2019darkgravity,cryptonote2012whitepaper}.
\item \textbf{A two-layer architecture.} A conventional UTXO Layer 1 for durable
settlement, and a Layer-2 task economy that is replayed deterministically from
Layer-1 blocks rather than run by a separate consensus \cite{dwork2012mining,
trivedi2017sok}.
\item \textbf{A trust-minimal bridge.} Layer-2 deposits and exits are authorised
by on-chain inclusion proofs rather than by a privileged operator, and the exit
crank is permissionless \cite{wood2014decentralized}.
\item \textbf{An agent-first authorisation layer, enforced by the chain.} A grant
published on chain binds its owner to a named spender and a named agent, and
the chain enforces the per-transaction cap, the tumbling per-window cap, the
allow-lists, the validity window and the nonce; a delegated grant that widens
its parent is refused, and revocation reaches every descendant. The same
applies to credit: the lender signs, the principal moves, the collateral
answers for the debt \cite{yao2023survey}.
\item \textbf{A second implementation of those rules.} The authority and credit
rules are written twice --- once in the protocol, once in an independent
reference derived from the specification --- and the two are compared scenario
by scenario, so a rule that refuses for a different reason is a defect rather
than a matter of taste.
\item \textbf{A documented security programme.} An explicit threat model, a
concrete list of implemented mechanisms, an audit-driven remediation history, a
fuzzing campaign, and an honest statement of the gaps that remain.
\item \textbf{An open governance posture.} Permissive dual licensing, community
ownership of the name, separation of the trademark from the licence, and a stated
refusal of post-launch governance in the style of Bitcoin
\cite{nakamoto2008bitcoin,bonneau2018sok}.
\end{itemize}

\subsection{Scope and honesty statement}

This paper describes the protocol as implemented in the reference repository
\cite{jounit-repo,jounit-spec}. The chain is \emph{pre-launch}. The repository's
own readiness record states plainly that no independent third-party security
review has been completed, and this paper does not claim otherwise
\cite{jounit-readiness}. Where a feature is experimental, that is said in the
text. Numbers quoted as measurements --- test counts, fuzzing executions,
benchmarks --- are reported as they appear in the repository's verification
record \cite{jounit-status,jounit-fuzz}, because the project has a deliberate
policy of replacing fabricated measurements with either a measurement or an
explicit ``not measured''.


\section{Jounit at a glance}
\label{sec:glance}

Jounit is a cryptocurrency with two layers that have different jobs.

\paragraph{Layer 1} is a conventional UTXO chain: wallets, transfers, blocks,
a peer-to-peer network, a fee market and a bounded issuance schedule. Nothing
about Layer 1 is specific to artificial intelligence. It exists to be a durable
settlement layer, and it deliberately inherits the well-understood properties
of Bitcoin \cite{nakamoto2008bitcoin,antonopoulos2017mastering} rather than
inventing a new ledger model.

\paragraph{Layer 2} is the part that makes the compute work worth doing. It is a
task economy. A \emph{submitter} posts an escrow for a described computation; a
\emph{worker} claims it, executes it and submits a result hash; anyone may
\emph{challenge} that result; a committee then decides whether the work was
correct, and settlement moves the money accordingly.

\paragraph{Layer 2 also carries the authority and the credit.} A \emph{grant}
is the owner's signed permission for a named agent to spend the owner's money
within stated limits; a \emph{position} is a loan one agent made to another.
Both are consensus state, committed by the block's Layer-2 state root, so every
node enforces the same limits rather than each node deciding for itself. This
is the part of the design that makes the paper's central claim usable rather
than merely stated: a fleet of agents can be given bounded spending authority
that the network itself checks, and one agent can lend to another with a
signature, a collateral and a deadline that every node can see.

\begin{figure}[t]
\centering
\fbox{\parbox{0.93\textwidth}{
\centering
\textbf{Layer 2 --- authority and credit}\\[2pt]
\textit{grant $\rightarrow$ spend $\rightarrow$ revoke; loan $\rightarrow$
collateral $\rightarrow$ repay $\rightarrow$ default}\\[6pt]
\rule{0.9\textwidth}{0.4pt}\\[6pt]
\textbf{Layer 2 --- task economy}\\[2pt]
\textit{post $\rightarrow$ claim $\rightarrow$ submit $\rightarrow$
(dispute?) $\rightarrow$ settle $\rightarrow$ exit}\\[6pt]
\rule{0.9\textwidth}{0.4pt}\\[6pt]
\textbf{Layer 1 --- UTXO settlement}\\[2pt]
\textit{wallets $\rightarrow$ transactions $\rightarrow$ blocks $\rightarrow$
hybrid proof of work $\rightarrow$ fees, burn, issuance}\\[6pt]
\rule{0.9\textwidth}{0.4pt}\\[6pt]
\textbf{Peer-to-peer layer}\\[2pt]
\textit{framing, handshake, relay, synchronisation, bans}
}}
\caption{The three-layer structure of Jounit. Layer-2 transactions are carried
inside Layer-1 blocks and committed by the block's Merkle root; Layer-2 state is
rebuilt by deterministic replay, so no second consensus protocol exists.}
\label{fig:layers}
\end{figure}

\subsection{Design principles}

Five principles recur throughout the codebase, and they explain most of the
otherwise surprising choices.

\paragraph{One code path, three profiles.} The development specification
requires that the fast test profile be \emph{the same code} with different
parameters, never a second implementation. Every consensus constant therefore
lives in one parameters module, and the \texttt{mainnet}, \texttt{testnet} and
\texttt{mini} profiles differ only in numbers. The practical consequence is that
anything tested on \texttt{mini} is the production code path, and that a
consensus bug cannot hide in a profile-specific branch.

\paragraph{Solving is hard, verifying is easy.} This is the property that makes
a blockchain a blockchain \cite{nakamoto2008bitcoin}. Jounit keeps the property
and changes the work: the proof of work is not pure hashing, so that hardware
useful for AI workloads is not left idle, but the security-critical part of the
puzzle remains a plain SHA-256 chain.

\paragraph{No floating point, anywhere.} The consensus core computes in integers
end to end. Weights derive from the epoch identifier by SHA-256, forward passes
accumulate in \texttt{i64} and wrap to \texttt{i32}, and every arithmetic
operation that could overflow is refused rather than wrapped. A wrapped sum
would compare small and pass a bound --- that is a security property, not a
style preference.

\paragraph{Fail closed.} Ambiguous cases refuse rather than guess. An empty
recipient list in a spending policy authorises nothing (not ``anything''); a
node with no configured credit issuer refuses every credit mint; a checkpoint
presented to a node with no committee fails rather than verifying against
nothing.

\paragraph{Structural, not documentary.} Where a guarantee matters, it is
expressed in the type system or in the shape of a data structure rather than in
a comment that a future maintainer could forget. The canonical reference that
decides a fraud dispute is \emph{computed by the arbitration engine} from the
task's committed fields and cannot be passed in as an argument. The credit
module exposes no function that accepts a reputation score at all, so the
guarantee is enforced by the compiler.


\section{Provenance: who built it, and how}
\label{sec:provenance}

\subsection{Authorship and ownership}

Jounit was conceived, specified, implemented and documented by a single author,
Swamy Gadila, and is published under a permissive dual licence
(\texttt{MIT OR Apache-2.0}). The author has disclaimed individual ownership of
the project: the code and the name are community property in the Bitcoin sense
\cite{jounit-repo,jounit-readiness}. The trademark treatment is deliberately
separate from the licence --- the name and ticker are community marks that
anyone may use to refer to the project factually, but a fork must take a
different name and a different ticker, so that two chains cannot both display
\texttt{JNT}. Section~\ref{sec:governance} discusses this in full.

\subsection{How the protocol was developed}

The development method is itself part of the result, because the protocol's
reliability rests on it more than on any single cryptographic choice. Three
phases ran in parallel from the beginning.

\paragraph{Specification first.} A single normative document (\texttt{SPEC.md},
roughly 2400 lines across 21 top-level sections) defines every wire format,
every consensus rule, every fee and every domain separator, together with
hand-written test vectors. The project's operating rule is that where a
disagreement exists between the specification and the code, the specification
wins --- which is only meaningful because the specification is written to be
executable in prose.

\paragraph{A decisions log.} A parallel document records numbered decisions
(D1--D24, plus P-series and M-series entries) with the reasoning behind each
one, including the decisions that were reversed and why. Several of the most
important properties of the finished protocol are direct consequences of
decisions recorded here: a rate-anchored issuance rule instead of a stock cap,
a stake-derived validator committee instead of a declared one, a proportional
settlement fee instead of a flat one, and a post-quantum succession rule that
is age-gated so that history is never stranded \cite{jounit-decisions}.

\paragraph{Independent re-implementation.} A second implementation of the
consensus-critical primitives --- the proof-of-work digests, the wire format,
the supply schedule and the difficulty rule --- was written in Python directly
\emph{from the specification}, not from the Rust source, and is used to
cross-check the Rust implementation. Agreement between two implementations
written from two different descriptions of the same rules is evidence that the
rules are unambiguous rather than merely self-consistent.

\subsection{The remediation culture}

The repository's verification record \cite{jounit-status} documents the
occasions on which a claim did not survive measurement. An earlier benchmark
document was deleted because it attributed numbers to a machine that no probe
had produced; an earlier fuzzing table was deleted as fabricated and replaced
with a measured campaign; eight mismatches between the specification and the
code were found and corrected. A protocol that is pre-launch can afford this
kind of honesty, and a protocol that is post-launch cannot. The paper therefore
treats the audit history (Section~\ref{sec:security-process}) as a first-class
part of the result.


\section{Architecture and implementation}
\label{sec:architecture}

\subsection{Crate layout}

The implementation is a single Rust workspace with seven crates
\cite{jounit-repo}.

\begin{table}[t]
\centering
\begin{tabular}{p{0.15\textwidth}p{0.13\textwidth}p{0.62\textwidth}}
\toprule
\textbf{Crate} & \textbf{Role} & \textbf{Notable contents} \\
\midrule
\texttt{jnt-core} & Consensus & chain, blocks, ledger, PoW puzzle, difficulty,
  supply, fees, bonds, finality, BLS, signatures, escrow, credit, L1/L2 bridge \\
\texttt{jnt-l2} & Layer 2 & task and service contracts, disputes, deposits,
  exits, verified work certificates, replay \\
\texttt{jnt-net} & Networking & framing, handshake, messages, relay, sync,
  address book, bans \\
\texttt{jnt-store} & Persistence & append-only checksummed block log, metadata
  file, node-local state container \\
\texttt{jnt-node} & Node & CLI, daemon, REST API, wallet file, mediator \\
\texttt{jnt-worker} & Worker SDK & discovery, claim, execute, submit, dispute,
  exit; client only \\
\texttt{jnt-mcp} & Agent surface & keyless dispatcher over identity, payments,
  receipts \\
\bottomrule
\end{tabular}
\caption{Workspace layout. All consensus logic lives in \texttt{jnt-core}.}
\label{tab:crates}
\end{table}

The architectural invariant that shapes everything else is that
\texttt{jnt-core} contains \emph{all} consensus logic and has no I/O, no clock
and no network access. It is a pure state machine: given a block and the
current state, it returns a new state or an error. This is what allows the
Python cross-checks, the fuzz harness and the property-based test suites to
exercise the real consensus path rather than a mock of it. The same discipline
governs the policy layer, which is documented as pure and deterministic so that
a wallet and a service that reach the same inputs reach the same answer.

\subsection{Concurrency model}

The node deliberately avoids an asynchronous runtime. A single thread owns the
chain; peer threads exchange messages with it through channels. This is not an
oversight: mining is an unbounded CPU loop, and sharing the chain under a mutex
across a thread pool invites lock contention exactly when the node is busiest.
The release profile additionally sets \texttt{lto = "thin"},
\texttt{codegen-units = 1}, \texttt{panic = "abort"} and, importantly,
\texttt{overflow-checks = true}, so an arithmetic overflow aborts the process
rather than wrapping silently. The protocol crates contain no \texttt{unsafe}
blocks at all.

\subsection{Storage and crash semantics}

Two storage artefacts exist and are deliberately kept apart.

\paragraph{The chain store} is an append-only log of checksummed blocks plus an
atomically replaced metadata file. The crash semantics are explicit: a torn
tail is dropped \emph{and reported}, whereas damage to committed data causes
the store to refuse to open rather than to silently truncate history.

\paragraph{The node-local books} are written to a separate file,
\texttt{state.bin}, using a simple container format:
a magic value, a version, a length, a SHA-256 digest of the payload, and the
payload itself. The file is written atomically and loaded fail-closed. The
rationale is stated in the source: a node that cannot read its own books must
refuse to start, because ``I forgot every credit line'' and ``I have no credit
lines'' are the same bytes and opposite facts. The envelope version has advanced
four times as features were added --- settlement links, escrow-to-exit
bindings, and finally the credit-attestation replay guard --- and each bump is
visible in the revision history.

\subsection{Continuous integration}

The continuous-integration pipeline enforces, on every change: formatting,
\texttt{clippy} with warnings escalated to errors, the full workspace test suite,
a release build, cross-checks of the consensus vectors against the Python
reference, the attack simulations, and a short fuzzing pass over every
registered target. A separate weekly workflow runs dependency auditing and
licence-policy checks. All third-party actions are pinned to commit hashes, and
automatic dependency updates are configured to ignore moving references and
semver-major bumps of the cryptography crates specifically, so that a routine
dependency bump cannot silently change a hash function or a signature
implementation.


\section{The blockchain layer}
\label{sec:l1}

\subsection{Encoding and identity}

The block header is exactly 177 bytes and contains a version, a height, a
64-byte previous-hash pointer, a Merkle root, a millisecond timestamp, a base
fee, a 128-bit difficulty, a nonce, a puzzle identifier, an epoch identifier, a
16-byte extra nonce and a one-byte signature-scheme tag. Transaction outputs
are 28 bytes: an 8-byte little-endian value and a 20-byte address.

Hashing is split by purpose. Internal commitments, Merkle nodes and proof-of-work
evidence use SHA-256 \cite{nist-fips180-4}. The two values an attacker most
wants to forge --- a transaction identifier and a block identifier --- use
SHA-512, on the explicit grounds that a 256-bit digest retains only about 128
bits of preimage resistance against a Grover-equipped adversary and about 85
bits of collision resistance against a quantum collision search, while doubling
the output width restores a full 256-bit classical and quantum margin
\cite{shor1994algorithms,grover1996fast}. Every hash input is domain-separated:
distinct tags separate block hashing, transaction hashing, transaction-body
signing, address derivation and each stage of the proof of work, so that a
digest computed in one context can never be replayed into another.

The Merkle tree separates its leaves by type as well. Layer-1 leaves, Layer-2
leaves, the Layer-2 state root and exit claims each carry distinct domain
tags, so an attacker cannot present a Layer-2 payload as a Layer-1 transaction
or lift a state root into a transaction position.

\subsection{The hybrid proof of work}
\label{sec:pow}

The proof of work is a three-stage puzzle, specified in
Section~5 of the protocol specification and implemented in the \texttt{puzzle}
module \cite{jounit-spec}.

\paragraph{Stage 1 --- the SHA-256 chain.} The header is serialised and hashed
with a domain tag and a 64-bit nonce; the result is then iterated for a fixed
number of rounds, each round hashing the previous value together with a
round index. This stage carries the security of the system. It is the part a
quantum adversary must break.

\paragraph{Stage 2 --- the integer-only neural network.} A small feed-forward
network with three layers of sixteen hidden units is evaluated on the Stage-1
output. The weights are not stored anywhere: they are derived from the epoch
identifier by hashing, which means no miner can precompute them before an epoch
begins, and every miner in the same epoch computes the same function. The
accumulator is a 64-bit signed integer that wraps to 32 bits at each layer, and
the activation is a shift-and-clamp. The bound on the accumulator is derived in
the specification and holds by construction, so no floating-point rounding
difference can arise between implementations.

The neural stage is a \emph{cost model}, not the load-bearing part of the
proof, and two mechanisms enforce that. The first is a work-share band: the
registry entry must declare that the neural stage accounts for between 30\% and
50\% of modelled work, and an entry outside that band is rejected outright, so
the neural stage can never become the weak link it exists to make useful. The
second is the ordering of the stages: the security-critical hash is Stage 1,
and Stage 3 binds the Stage-1 and Stage-2 outputs together with the nonce and
epoch identifier into the single value that must fall below the target.

\paragraph{Stage 3 --- binding.} The final digest hashes the domain tag, the
epoch seed, a digest of the activations, the nonce and the epoch identifier.
Validity is the usual comparison against a target derived from the block's
declared difficulty:
\begin{equation}
\mathrm{U256BE}\big(\mathrm{SHA256}(D_{S3} \parallel \mathrm{seed} \parallel
d_{\mathrm{act}} \parallel \mathrm{nonce} \parallel \mathrm{epoch})\big)
\leq \left\lfloor \frac{2^{256}-1}{d} \right\rfloor ,
\label{eq:pow}
\end{equation}
where $d$ is the block difficulty.

A second proof-of-work family, a memory-hard scratchpad puzzle, has been
specified and implemented and has passed a scoped internal security review
\cite{jounit-hbm-review}, but it is deliberately \emph{not registered} as an
active puzzle identifier: an unregistered puzzle validates nothing, so the
alternative cannot silently displace the deployed one.

\subsection{Difficulty adjustment}

Difficulty follows a windowed, linear-weighted moving average in the style
established by the CryptoNote work \cite{cryptonote2012whitepaper}, with a
window of sixty blocks and a maximum per-adjustment change of ten per cent. The
retarget is multiplicative rather than additive:
\begin{equation}
d_{\mathrm{new}} = d_{\mathrm{prev}} \cdot \frac{T}{\bar{t}},
\label{eq:retarget}
\end{equation}
with $T$ the target block time and $\bar{t}$ the observed average solve time.
Additive adjustment would bias the schedule in one direction under sustained
pressure, and multiplicative adjustment has the additional property that the
relative change is bounded by construction.

Timestamps are checked against a lower bound only: the median of the previous
eleven blocks' timestamps. There is no consensus upper bound, because a
clock rule in the consensus path is a lever that a miner can pull to influence
difficulty; instead the node applies a local sanity limit one block time beyond
the target.

\subsection{Fork choice}

The chain is linear and hash-linked; there is no DAG and no uncle structure.
Fork choice is heaviest cumulative work: the node keeps the branch with the
greatest accumulated difficulty, with ties broken deterministically by lower
height and then by the lowest block hash, so that two nodes holding the same
set of valid blocks always agree on the tip. Cumulative work is computed
incrementally at insertion time by saturating addition, and side branches are
pruned once they exceed a threshold.

Because the fork-choice rule is a total order over the set of valid blocks,
the system inherits the security analysis of cumulative-work chains
\cite{eyal2014securitybitcoin,kiayias2017equitable}: an attacker must spend
real work to displace the honest branch, and the cost of a reorganisation
grows with the work it must discard.

\subsection{Fees, burn and the protocol reserve}

The fee market is a base-fee mechanism in the style of \texttt{EIP-1559}
\cite{buterin2020eip1559}: a required fee proportional to gas, a base fee that
rises when blocks are full and decays otherwise, clamped to a configured floor
and ceiling. Gas is charged per transaction, per input and per output.

The base fee is split in half: one half is burned, the other half is credited
to an on-chain reserve used for protocol security and compute budgeting. The
split is executed with integer arithmetic that assigns the odd satoshi to the
burned side, on the reasoning that a burn is irrecoverable and conservative
rounding in that direction is the safe choice. Only the tip above the base fee
is paid to the miner.

\subsection{Issuance}
\label{sec:issuance}

There is no stock cap. Issuance is rate-anchored and time-based
\cite{jounit-monetary}, and three properties of the schedule are worth stating
because they are unusual and because they are enforced in code rather than in
prose.

\paragraph{Zero premine.} The genesis block mints nothing. This is enforced by
the code path, not merely by the absence of a special case.

\paragraph{A halving ladder in whole units.} The block reward halves every
262\,800 blocks --- approximately five years at a ten-minute block time --- in
whole units of the token rather than in shifted subunits, because halving
$9.5$ to $4.75$ breaks the arithmetic. The ladder is
\begin{equation}
76,\ 38,\ 19,\ 9,\ 4,\ 2,\ 1 \ \ \mathrm{JNT},
\label{eq:ladder}
\end{equation}
and the total issuance over the ladder era is exactly 39\,157\,124 JNT. That
figure is not the geometric shortcut: the genesis block occupies a rung-zero
height without minting, so the first rung is paid by one fewer block than a
naive computation suggests, and the difference is precisely the shortfall
between the approximation and the chain's real reach.

\paragraph{A tail that never stops.} After the ladder is exhausted, the reward
decays to a floor rather than to zero: $0.50$, $0.25$, $0.125$ JNT, and then
a permanent floor of $0.1$ JNT per block, forever. Issuance therefore never
terminates, and net supply can actually fall, because the burned half of the
base fee removes supply from the spendable set. The protocol's own monetary
analysis reports that net supply crosses below break-even at roughly five per
cent activity and sits at about a quarter of break-even at the fee floor. This
is a measured consequence of a deliberate design, not an artefact.

Two bounds enforce the schedule independently of the miner, the validator and
the node, all checked when a block is applied. The first is an issuance
integral: a block may not mint past the cumulative issuance at its height, and
overflow is refused rather than wrapped. The second is a rate anchor: after
ladder exhaustion, gross issuance inside the current halving window may not
exceed one per cent of the supply issued at that window's start, which bounds
dilution at roughly 0.2 per cent per year in the limit. The approved schedule
sits about three times under the anchor, so it never binds today; it exists to
refuse a \emph{future} schedule that would dilute far faster.

\subsection{Finality}
\label{sec:finality}

Jounit uses two layers of finality, and the distinction is important because
the two protect against different events.

\paragraph{Depth-based finality} is the familiar rule: a transaction buried
under a configured number of blocks is considered settled. On the mainnet
profile this is one hundred blocks. It is cheap and probabilistic.

\paragraph{Binding finality} is a quorum certificate over a block height,
signed by a stake-weighted committee using BLS signatures
\cite{boneh2003aggregate,bls-signature2022}. A checkpoint is accepted only if
it satisfies a seven-clause checklist: epoch alignment, a self-derived signer
set, an exactly-sized signer bitmap with no bits outside the set, a popcount
equal to the declared signer count, satisfaction of a quorum threshold,
verification of the BLS aggregate signature, and the presence of a
post-quantum sidecar where one is required. A certificate at height $H$ binds
the chain at $H$ only if the canonical block at $H$ exists, the certificate
validates, the quorum is met, the height is strictly greater than the previous
binding point, and no conflicting certificate exists for the same height.

The quorum threshold is defined as a ceiling of two thirds:
\begin{equation}
q(n) = \left\lceil \frac{n \cdot 6667}{10000} \right\rceil ,
\label{eq:quorum}
\end{equation}
giving three of four and fifteen of twenty-one for the relevant committee
sizes. The committee is derived locally from the bond ledger by stake-weighted
selection --- never from a list a peer supplies --- so a message is always
recorded against the receiving node's own view of the committee and height,
rather than against the sender's say-so.

Equivocation is detected structurally: a conflicting pair of certificates from
the same signer is admissible evidence, is bounded in size, and can be
submitted permissionlessly. Reorganisations that would cross the binding
finalised point are refused outright, so a finalised history cannot be undone.

\subsection{Peer-to-peer layer}

The transport is TCP, with one thread per peer and no asynchronous runtime.
Framing is \texttt{magic $\parallel$ type $\parallel$ length $\parallel$
payload}; the magic value is re-checked on every frame rather than only at
connect, the payload length is bounded before any allocation is made, and an
oversize frame closes the connection rather than being skipped, because a
stream protocol has no resynchronisation point. Payloads are read incrementally
so that a lying length field pins only the bytes actually delivered.

The handshake exchanges a network magic, a network identifier, a protocol
version, a service bitmask, the peer's height and tip hash, a user agent and a
nonce, and the responder must echo the initiator's nonce. This is
\emph{transport-level} peer identity only: there is no node key and no
authentication in the current version, and the protocol specification says so
rather than implying more.

Relay is bounded on every axis that a flood could exploit: the number of
inventory items per message, the number of transactions requested per batch,
the number of in-flight and pending relay items, the number of known
transactions remembered, and separate per-peer per-minute budgets for
transaction relay and for control messages. A node never relays a transaction
it did not itself admit, and never relays a block it has not validated.

Misbehaviour is scored rather than instantly banned: a ban score accumulates,
decays at one point per hour, and a threshold bans the peer. A block that fails
validation costs a small score increment rather than an outright ban, because
a peer relaying one bad block is not the same as a peer attacking the node.
Connection counts are capped (125 peers, of which at most 32 inbound), and
in-flight handshakes are counted separately, since a handshake that has begun
but not completed is invisible to the peer count and would otherwise let a
single host pin one thread and one socket per connection.

Finally, there is no peer discovery protocol at all: no address gossip, no
DNS seeds, no super-nodes. Bootnodes and a local peer file are the only peer
sources. This is a scoping decision, recorded as such, and it means that the
network's resistance to partitioning rests on the operator rather than on the
protocol --- a limitation we prefer to state rather than hide.


\section{The task economy}
\label{sec:l2}

\subsection{Design}

Layer 2 is not a second consensus protocol. Layer-2 transactions are carried
inside Layer-1 blocks and committed by the block's Merkle root, and Layer-2
state is rebuilt from that chain by deterministic replay. What does not exist
is any rule that makes a Layer-1 block invalid because of its Layer-2 content:
a Layer-2 payload that fails to apply is recorded by the replayer and skipped
deterministically, never converted into a consensus rule. This is a
deliberate separation --- Layer 2 inherits Layer 1's ordering and immutability
without inheriting Layer 1's latency, and the failure mode of a malformed
Layer-2 payload is a skipped payload rather than a halted chain.

The vocabulary of the task economy is six transaction types.

\begin{table}[t]
\centering
\begin{tabular}{cll}
\toprule
\textbf{Code} & \textbf{Type} & \textbf{Purpose} \\
\midrule
\texttt{0x10} & Post & Submitter opens a task, locks the escrow, publishes the
  compute specification \\
\texttt{0x11} & Claim & Worker takes the task and posts a deposit \\
\texttt{0x12} & Submit & Worker submits the result hash with a signed
  attestation \\
\texttt{0x13} & Dispute & Anyone challenges the result with a counter-result
  and a deposit \\
\texttt{0x14} & Settle & Committee applies the decision: pay, refund or slash \\
\texttt{0x15} & ExitQueue & Owner starts withdrawing value back to Layer 1 \\
\bottomrule
\end{tabular}
\caption{The six Layer-2 transaction types. A task moves through
$\text{Open} \rightarrow \text{InProgress} \rightarrow \text{ResultSubmitted}
\rightarrow \text{Disputed?} \rightarrow \text{Settled}$, with $\text{Expired}$
as the timeout branch.}
\label{tab:l2types}
\end{table}

\subsection{Value conservation}

Every Layer-2 settlement enumerates the movement of every unit of value ---
worker paid, worker deposit returned, worker deposit slashed, submitter
refunded, challenger refunded, challenger deposit burned, worker deposit
burned, fee burned, fee to the reserve --- so that a settlement cannot be
under-specified in a way that creates or destroys value. The Layer-2 state
machine asserts a conservation invariant after every step:
\begin{equation}
\sum \text{accounts} + \sum \text{locked} + \text{burned} + \text{reserve}
+ \text{vault remaining} = \sum \text{credited} ,
\label{eq:conservation}
\end{equation}
and a fuzz target asserts exactly this property over adversarially generated
transaction sequences. The joint Layer-1/Layer-2 check needs no exit term at
all: the Layer-1 payment and the Layer-2 debit cancel while the sink term falls
by the same amount.

\subsection{Deposits and exits: the bridge between layers}
\label{sec:exits}

Deposits are not one-way, and the mechanism by which value returns to Layer 1
is the most security-sensitive part of the design.

\paragraph{Deposit.} A deposit is an ordinary Layer-1 payment to the profile's
derived sink address --- no new transaction type, no special opcode. The
amount is read from the committed output, never from the claim, and the
deposit identity is the pair (transaction identifier, output index). Credit
requires both depth-based finality and binding finality, so the rule for
maturity is never weaker than the rule used elsewhere in the protocol.

\paragraph{Queue.} The owner signs an exit-queue transaction. It debits the
owner by amount plus exit fee plus queue fee, and records
$(\text{owner}, \text{nonce}) \rightarrow (\text{amount}, \text{destination},
\text{exit fee}, \text{queued height})$. Nonces are per owner and strictly
increasing; zero amounts, under-floor fees and repeated or rewound nonces are
refused.

\paragraph{Maturity.} An exit becomes settleable after the same finality
requirement as a deposit --- depth and binding --- never less.

\paragraph{Release.} The claim is authorised by a Merkle inclusion proof
against the Layer-2 state root committed at the queue block, not by a
signature, because the sink holds no key. The verifier recomputes the exit
record leaf from the recorded fields and checks it against the committed root;
no caller-supplied amount, destination or fee is ever read from the claim.
The inputs consumed must sum to exactly the amount plus fees, with no change
and no tip, so an exit whose funding is not present stays queued rather than
draining a stranger's deposit. Consensus pays only the destination that the
proven record names, so a caller cannot redirect it, and the crank collects a
fee the owner prepaid at queue time --- which is what makes a crank nobody
wants to run a non-issue rather than a griefing vector. The crank is
permissionless: a read-only endpoint hands settleable claims to any third
party, so a non-miner can construct the block that includes them.

\paragraph{Reorganisation safety.} A release is an ordinary undoable Layer-1
operation. Rolling back the settling block restores the sink outputs, clears
the exit identity, and Layer-2 replay regenerates the watch. The exit identity
is consumed in consensus state, so a given $(\text{owner}, \text{nonce})$ pair
is released exactly once, ever.

\subsection{Disputes and adjudication}
\label{sec:disputes}

The interesting security question in a task market is: who decides, and can the
decider be influenced? Jounit's answer is structural.

\paragraph{The reference cannot be supplied.} The canonical reference result
that decides a fraud dispute is computed by the arbitration engine from the
task's own committed fields. It is never accepted as an argument. No caller can
dictate a verdict, because there is no parameter through which a verdict could
be dictated.

\paragraph{Evidence and verdicts.} A dispute carries evidence --- a
counter-result and the associated attestations --- and the engine produces
exactly one of three verdicts: the worker was correct, the challenger was
correct, or the dispute is inconclusive. Inconclusive is deliberately not a
verdict: it settles nothing, slashes nothing and burns nothing, because a
system that must burn something when it cannot tell would punish honest
disagreement.

\paragraph{Consequences of a fraud verdict.} The worker's deposit is moved to
the vault, where half is burned immediately and half vests linearly over a
year, funding a whistleblower bounty of twenty-five per cent of the slashed
amount paid out of the reserve. The submitter is refunded in full; the
challenger is refunded and receives the bounty. No certificate is issued for a
fraudulent or expired task.

\paragraph{Committee voting.} Votes are transported over the peer-to-peer
network as arbitration votes and tallied against a bounded book. Every vote is
checked against the receiving node's own derived committee, and votes are
recorded for the height they concern rather than for the height of the
sender's claim. The committee meets a quorum of two thirds for a decision to
bind.

\subsection{Verified Work Certificates}
\label{sec:vwc}

A verified work certificate is the protocol's receipt for work that settled
as correct. It records that task $T$ settled correct with result $R$, worker
$W$ was paid $P$ at height $H$.

It is a \emph{derived consequence} of a settlement, not a claim an actor
makes. There is no issuance transaction and no issuer key, because the
settlement path already performed all the authorising. Its identifier is
derived from every other field, so a tampered certificate cannot carry a
matching identifier, and its encoding is byte-unique.

Two kinds exist, and the difference between them is load-bearing:

\begin{table}[t]
\centering
\begin{tabular}{@{}p{0.16\textwidth}p{0.55\textwidth}p{0.22\textwidth}@{}}
\toprule
\textbf{Kind} & \textbf{What backs it} & \textbf{Offline-verifiable?} \\
\midrule
\texttt{Quorum} & The settle payload's own BLS aggregate over the decision &
  Yes, against the validator set \\
\texttt{WindowElapsed} & Nothing beyond the on-chain record that nobody
  objected in time & No \\
\bottomrule
\end{tabular}
\caption{Two kinds of verified work certificate. The weak basis is unmissable
rather than a convention a relying party could forget: the API refuses the
window-elapsed kind by name when a correctness claim is required.}
\label{tab:vwc}
\end{table}

The unchallenged-result rule deserves emphasis because it is the honest weak
point of any optimistic settlement system. If a worker submits a result and
nobody challenges it within the dispute window, the task auto-settles correct
and the worker is paid. That is correct protocol behaviour: a challenge must be
able to overturn a settlement, because otherwise anyone could file a dispute
and never resolve it, stalling every task forever. The consequence is that an
incorrect result that nobody challenges in time \emph{is} paid. What must never
happen is that such a payment being mistaken for verification. Three
mechanisms prevent it: the certificate kind is explicit in the data structure,
the validation function refuses the weak kind by name, and the certificate
assertion returns false for it --- so a relying party that needs independent
proof must demand the quorum basis, and there is no argument by which the weak
one satisfies it.

\subsection{Reputation}

Reputation is a bounded accumulator per identity, driven by a fixed weight
table held in code rather than in data: service completed $+1$, payment
completed $+1$, verified work $+2$, dispute won $+1$, dispute lost $-3$. An
event carries no amount, no weight and no signature of its own, so a
differently-shaped event cannot forge a better score; and losing a dispute
costs three times what winning one earns, so baseless challenges are
expensive.

Critically, reputation is \emph{not} a consensus input. A credit limit computed
on-chain from a mutable score would make every reputation update a consensus
event and make fresh identities farmable, which is the cheapest available
attack against the gate. The credit module therefore exposes no function that
accepts a score at all: the guarantee is the type system refusing the call,
not the code declining it.


\section{Agent authority and credit on chain}
\label{sec:authority}

\subsection{Why it had to move}

The spending-policy layer of an earlier revision of this protocol was
\emph{node-local}. A policy constrained a payment on the node that issued it,
and the chain never saw it: the consensus core contained no reference to a
policy, an agent, a delegation or a credit at all. A machine fleet bound by
policies was therefore bound by the goodwill of whichever node had issued them,
and two nodes holding the same chain could disagree about whether a spend was
permitted.

The credit facility had the same shape and two further problems. Depositing
credited a local balance without moving a coin, and opening a position accepted
a \emph{lender} field that anybody could write: a loan the lender had never
agreed to, on money that did not exist, visible only to the node that recorded
it.

Both are now Layer-2 consensus state. The distinction that matters is not
where the code lives but what a node is \emph{required} to do with it: a rule
that every node replays is a rule, and a rule one node checks is advice.

\subsection{The grant}

A grant is the owner's signed permission. It is published with a Layer-2
transaction, and the owner's signature covers three things at once --- the grant
bytes, the \emph{spender} (the key that may exercise the grant) and the
\emph{agent identity} the grant names. That binding is the whole point: a grant
hash on its own is a capability, and a capability anybody can exercise is
worthless. Here the owner says \emph{this key may spend under these limits},
and the chain remembers it.

The grant carries the limits the owner stated: a per-transaction cap, a
tumbling per-window cap, an exhaustive recipient allow-list, an exhaustive
service allow-list, a validity window, and a parent link for a delegated grant.
The chain enforces all of them against state it keeps:

\begin{enumerate}
\item the per-transaction cap: $a \leq C_{\mathrm{tx}}$;
\item the per-window cap, over a tumbling window that starts at the first spend
  and rolls over after $W$ blocks:
  \begin{equation}
  \sum_{\text{spends in the window}} a \;\leq\; C_{\mathrm{window}} ;
  \end{equation}
\item the recipient allow-list, where an \emph{empty} list allows nothing;
\item the service allow-list, where a grant that names services authorises no
  untagged spending;
\item the validity window, which is inclusive at its start and exclusive at its
  end;
\item a strictly increasing nonce, which is what makes an authorisation
  unreplayable;
\item and the rule that a height inside a grant never runs backwards.
\end{enumerate}

One grant is live per agent. A second publication for the same agent is
refused, so a second grant cannot be used to double the caps the first one
grants. A delegated grant must \emph{narrow} its parent --- the child's caps are
not wider, its allow-lists are not wider, and it does not outlive it --- and the
narrowing is checked against the parent's bytes as they are held on chain, not
against a claim about them.

\subsection{Revocation, and what it reaches}

The owner can kill a grant, and the kill switch reaches every descendant: a
spend checks the whole ancestry, so a delegated fleet dies with the grant it
descends from, without a revocation transaction per link. The walk is bounded
by the delegation depth cap and fails closed --- a missing ancestor, a revoked
one, or a chain deeper than the cap all mean \emph{not live}.

\subsection{Spending}

A spend is signed by the spender, not by the owner. The owner signed once, at
publication; the grant itself authorises every later spend. The money comes
out of the grant owner's balance --- a grant is an allowance over the owner's
own funds --- and the fee comes out of the same balance on top.

\subsection{Credit}

A loan is four transaction types. The \emph{lender} signs the terms; the
principal moves to the borrower's balance when the position opens; the
borrower's own funds are locked as collateral, and they are checked
\emph{before} the principal arrives, so the borrowed money cannot secure
itself. Interest is simple interest per whole window on the outstanding
principal,
\begin{equation}
\text{interest} \;=\; \frac{\text{principal} \times r_{\mathrm{bps}} \times
\text{windows}}{10\,000} ,
\end{equation}
charged once per period rather than once per repayment, and a repayment pays the
interest before the principal. A repayment larger than the debt is refused: the
lender is paid the debt and not a satoshi more.

A default is permissionless --- anyone may trigger the settlement of an overdue
position, because a default only the lender could trigger is a default that
could be delayed. The collateral answers for the debt, the surplus returns to
the borrower, and what the collateral does not cover is a shortfall the lender
absorbs rather than money the protocol invents.

\subsection{The state commitment, and the invariant}

Both families live in the Layer-2 state, under their own leaf tags, and are
committed by the block's Layer-2 state root. The conservation invariant that
already covered the task economy now covers them too:
\begin{equation}
\sum \text{accounts} + \sum \text{locked} + \sum \text{exits locked}
+ \sum \text{collateral locked} + \text{burned} + \text{reserve}
+ \text{vault remaining} \;=\; \sum \text{credited} ,
\end{equation}
asserted after every single transition, including every refusal.

\subsection{How this is checked}

The rules of this section are implemented twice: once in the protocol, and once
in an independent reference written from the specification in a different
language. The reference derives the addresses, the agent identities, the grant
hashes and the position identifiers itself, and emits a vector file of
scenarios; the protocol's test suite replays those scenarios and asserts the
same verdict, the same reason and the same final state. A rule that refuses for
a different reason is a different rule, and the two implementations have to
agree on which reason it was.

\section{Value, wallets and market operations}
\label{sec:wallets}

\subsection{Keys and addresses}

Jounit supports three wallet modes --- classical, hybrid and post-quantum ---
and the address is derived from the \emph{whole key set}, not from a single
key. A hybrid bundle contains an ECDSA component and a Schnorr component; the
post-quantum bundle adds an ML-DSA-65 component \cite{nist-fips204}. Because
the address commits to the serialised bundle, a hybrid key set cannot quietly
be replaced with a smaller one, and the three modes produce three distinct
addresses by construction.

Two rules govern signature bundles and both are enforced in code: every
component must verify, and the address commits to the whole key set. The
second is what makes downgrade attacks against a hybrid signature impossible.

The post-quantum requirement is age-gated. It is off on every current profile,
and when it is switched on it applies only to outputs created at or after the
activation height, so that historical outputs are never stranded. This is a
deliberate reading of the general transition problem: the cost of migrating
history is paid by nobody in particular, and the benefit of an early switch
does not justify stranding coins \cite{nist-ir8547}.

\subsection{Key management}

Wallet keys are held in a file whose format is a magic value, a random salt, a
random nonce and a ChaCha20-Poly1305 ciphertext \cite{rfc8439}, with the
key-encryption key derived from an operator passphrase by Argon2id
\cite{rfc9106}. The file is written atomically: a temporary file is created
with owner-only permissions \emph{before} any key material is written to it,
flushed, and renamed over the destination. A fresh salt and nonce are generated
on every write. The passphrase is read from an environment variable rather than
a command-line flag, because flags are visible in the process table. Storing
keys in plaintext requires an explicit opt-in flag; the default is refusal.

Secret hygiene is enforced at the type level as well: the wallet's private key
is wrapped in a zeroising type, the debug representation of a wallet redacts
the scalar, and the value is zeroised on drop.

\subsection{Storage and transfer of value}

Value is stored in an unspent transaction output model. Each output is a
28-byte structure; a transaction consumes previous outputs and creates new
ones; coinbase outputs mature after a configured number of blocks, and a
miner cannot claim two matured rewards in one block because outputs to the
same address accumulate rather than replace.

Transfers are ordinary: the node can construct and broadcast a transaction
locally or over the peer-to-peer network, and the CLI reports balances in whole
tokens while accepting amounts in subunits. A dedicated verification command
replays the entire chain from genesis and diffs the resulting output set,
supply, bond ledger and binding-finalised point against stored state --- the
check worth running after any change to the protocol.

\subsection{Buying and selling: the payment path}
\label{sec:payments}

There are two payment paths, and it is worth being precise about which is which.
The \emph{authority} path --- the one this paper's Section~\ref{sec:authority}
is about --- moves real value under a limit the chain enforces. The
\emph{service} path below is the operational bookkeeping around it: a request, an
authorisation, a receipt and a settlement link. The service path never invents
balance: a settlement still needs a transaction that actually pays.

For a machine-to-machine transaction, Jounit separates \emph{authorisation}
from \emph{payment}. The separation is what allows an agent to act without
holding a spending key in the general case.

\paragraph{Authorisation.} A principal signs a spending policy. The policy
answers exactly one question: may this agent, at this height, spend this much,
to this recipient? The implementation is pure and deterministic, and every
refusal is a named error rather than a boolean. The ambiguous cases fail
closed: an empty recipient list authorises nothing rather than anything; an
empty service list authorises no service-tagged payment; a nonce must strictly
increase, so an authorisation cannot be replayed; a height below the policy's
own window start refuses, so time does not run backwards inside a policy; and
cumulative spend is bounded per window, with the window rolling over rather
than accumulating forever. The same rules are the ones the chain enforces for
an on-chain grant, so an operator can reason about one set of limits rather
than two.

\paragraph{Attenuated delegation.} Delegation chains let an agent grant a
subset of its own authority downward. Each level inherits a strict subset of
its parent's policy --- never a wider one --- so authority narrows as it
travels. A payment under a chain names its own agent identifier, which is
derived from the ancestry it presents rather than from the payer's address;
the node records the links so that the lineage can be proved later, verifies
the owner's signature on every policy it grants under, and propagates a
revocation down to every recorded descendant. On chain, the same rule is
checked against the parent's bytes as the chain holds them
(Section~\ref{sec:authority}), so a delegated fleet's authority narrows for
every node, not only for the one that issued the grants. The practical result
is a fleet of agents whose operator can bound what any single agent can spend
without re-authorising each transaction by hand.

\paragraph{Payment.} A machine payment carries the authorisation, the
recipient, the amount and the reference. A receipt book records deterministic
service receipts and refuses a receipt that conflicts with an existing one for
the same payment. A payment tracker provides exact-payment replay protection
and refuses rather than evicts when its cap is reached.

\paragraph{Settlement.} A payment is not settled merely because it was
recorded. A settlement endpoint binds the recorded payment to a Layer-1
transaction that pays its recipient; a transaction that does not pay the
recipient is refused, a second, different transaction is refused as a
conflict, and the link persists across restarts. This is the mechanism that
prevents a node from claiming that a payment happened without the coins having
moved --- the same concern that motivates transaction inclusion proofs
elsewhere in cryptocurrency design.

\paragraph{Streaming micro-payments.} A payment stream is pay-as-you-use: a
payer-signed rate per block, a window, a hard cap, a recipient and the
governing policy for non-human payers. Settlement receipts are payer-signed,
the stream book is fail-closed (exact amounts, strictly ordered settlements,
no future-dated settlements, accrual checked for wraparound), and settlement
rides the ordinary payment path with the stream identifier as the reference, so
no new transaction type is introduced.

\subsection{Credit: borrowing against future work}
\label{sec:credit}

An agent with a demonstrated track record should not have to pre-fund every
transaction. The credit facility allows an agent to borrow against future
work. The loan itself is the on-chain facility of
Section~\ref{sec:authority}: the lender signs, the principal moves, the
collateral is the borrower's own money, and the interest and the default rules
are the chain's. What remains here is the service layer around it, and it is
worth being explicit that it is a \emph{gate}, not a source of spendable
balance.

\paragraph{Issuance is attested, not open.} New credit enters only through an
attestation signed by a key the operator listed as trusted. With no key
listed, every mint is refused --- the default is refusal, not implicit trust.
An attestation binds the network, the subject, the amount, an issue height, an
expiry height, the issuer key and the signature; it is checked against the
network, the height window, the trusted issuer list, the signature itself, and
a replay guard that records consumed attestations and persists across
restarts, so one attestation mints exactly once, ever.

\paragraph{Credit is always collateralised.} The required collateral is a
function of the amount, discounted by a posted bond but never replaced by it,
with a minimum collateral ratio enforced. On chain, where no bond is posted,
the requirement is full collateralisation. Settlement takes the position's own
collateral first, then the bond, and \emph{never} reputation. Collateral in
excess of the debt returns to the borrower: a default takes what it is owed
and no more. Any residual debt is a reported shortfall that the lender absorbs,
rather than the protocol inventing funds. A position runs to a deadline, and
before that deadline a default is refused --- a lender cannot settle a merely
late position on a whim --- and the bond must be posted before the position
opens, so no discount is ever unbacked.

\paragraph{One implementation of each rule.} The collateral arithmetic and the
interest arithmetic live in one place, and the on-chain facility and the
service-layer book both use them. Two implementations of the same rule would
be two rules.

\subsection{Conditional escrow}
\label{sec:escrow}

Conditional escrow is the general primitive behind the task economy, and it
also exists standalone for non-compute deliverables.

\paragraph{A contract-selected arbitrator.} Each escrow names its own
arbitrator as a BLS public key, bound when the escrow opens. There is no
committee, no central arbitrator and no on-chain registry: the contract
itself is the registry. The escrow identity is a hash over the network, the
parties and their keys, the arbitrator, the amount, a nonce and the condition,
so two escrows between the same parties with different terms are distinct
objects.

\paragraph{Release modes.} An arbiter release is verified as exactly one BLS
signature under the escrow's own arbitrator key; a mutual release requires
exactly two, and the claimed keys must equal the escrow's bound payer and
payee keys, so a release cannot be forged by a party that is not the party.

\paragraph{Expiry and refund.} Every escrow expires after a configured number
of blocks --- seven days on the mainnet profile --- and refunds its payer. A
refund requires no signature at all, deliberately, because it is the path for
``nobody will ever sign''; releases are refused after expiry, so a late
signature cannot race a refund.

\paragraph{Oracle conditions are refused.} Conditions that genuinely require an
oracle are refused when the escrow opens, because this codebase implements no
oracle verifier: an oracle-gated escrow would be a contract that nobody could
ever satisfy. Failing closed here is the difference between a mechanism that
is unavailable and a mechanism that is unsound.

\paragraph{Settlement to the chain.} A released escrow does not merely change
a book entry. The settlement endpoint binds the release to the Layer-1 exit
that pays the payee --- owner is the payer, destination is the payee, amount is
the escrow's amount --- and queues that exit through the ordinary withdrawal
path. The escrow reads finalised only once consensus actually released the
bound exit; a height stamp alone no longer finalises a bound escrow. This
means the escrow and the exit cannot disagree, because the escrow's final
state is defined by a Layer-1 fact rather than by a local one.

\subsection{Fee schedule}

Layer-2 fees use their own path and never consult the Layer-1 base fee, with
the flat fees shown below and a settlement fee proportional to the value
released (one tenth of one per cent, with a floor).

\begin{table}[t]
\centering
\begin{tabular}{lr}
\toprule
\textbf{Operation} & \textbf{Fee} \\
\midrule
Post task & 5\,000 subunits \\
Claim task & 2\,000 subunits \\
Submit result & 2\,000 subunits \\
Raise dispute & 10\,000 subunits \\
Settle & $\max(0.1\% \times \text{value},\ 1\,000)$ subunits \\
Queue exit & 2\,000 subunits \\
Exit fee floor & 1\,000 subunits \\
Layer-2 deposit floor & 1\,000\,000 subunits per sink output \\
\bottomrule
\end{tabular}
\caption{Layer-2 fee schedule. A proportional settlement fee means that a
proportional cost follows the value released, while a floor keeps small
settlements from being free to spam.}
\label{tab:fees}
\end{table}


\section{Security model}
\label{sec:security}

\subsection{Threat model}

The adversaries considered are the ones a public, permissionless, TCP-hosted
network actually faces: a remote attacker who wants to steal funds; a peer who
wants to exhaust a node's resources; a miner who wants to commit a block that
should be rejected; a block producer who wants to forge an exit claim; a worker
who wants to be paid for work it did not do; a submitter who wants a refund
without paying; a challenger who wants to grief a competitor; an API client
who wants to make the node misbehave or leak secrets; and a local attacker who
can read the node's data directory.

\subsection{What is trusted}

The trust model is best stated as a table, because the interesting entries are
the negative ones.

\begin{table}[t]
\centering
\begin{tabular}{@{}p{0.22\textwidth}p{0.11\textwidth}p{0.58\textwidth}@{}}
\toprule
\textbf{Component} & \textbf{Trusted?} & \textbf{Reason} \\
\midrule
Consensus core & Yes & The only trusted code; it has no I/O, clock or network \\
Peers & No & All rules are checked locally \\
Block producer & No & A lying Layer-2 root is accepted but only stalls exits \\
Exit claimant & No & The crank is permissionless and pays only the recorded
  destination \\
API caller & No & The API never signs and never submits to consensus \\
Model Context Protocol surface & No & It verifies with the same public
  functions; it cannot mint authority \\
Worker SDK & No & Every privileged transition is a signed Layer-2 transaction
  the protocol verifies \\
Reputation & No & Not a consensus input, by construction \\
Delegation service layer & No & The lineage book is the node's own record; the
  narrowing rule itself is a consensus rule (Section~\ref{sec:authority}) \\
Credit issuer key & Yes, locally & An explicitly configured operator-level
  trust anchor, and only for the service-layer book \\
Grant owner & Yes, by construction & A grant is the owner's own signature, and
  the chain enforces only what that signature says \\
\bottomrule
\end{tabular}
\caption{Trust boundaries. Note that the two ``yes'' entries are narrow: one is
the protocol's own code, the other is a key the operator chooses to trust and
can therefore choose not to trust at all.}
\label{tab:trust}
\end{table}

The absence of a trusted issuer by default is worth emphasising: credit minting
is fail-closed, so an operator who does not configure an issuer has a node that
refuses every credit request. The alternative --- an implicit issuer --- would
mean the operator had trusted something they never chose.

\subsection{Cryptographic mechanisms}

\begin{table}[t]
\centering
\begin{tabular}{@{}p{0.19\textwidth}p{0.27\textwidth}p{0.45\textwidth}@{}}
\toprule
\textbf{Primitive} & \textbf{Where} & \textbf{Purpose} \\
\midrule
SHA-256 & everywhere & commitments, Merkle nodes, proof-of-work stage 1, stage
  domain separation \\
SHA-512 & block and transaction identity & 256-bit classical and quantum
  margin on the two forgeable values \\
secp256k1 (ECDSA, Schnorr) & wallet and spending keys & transaction signing,
  hybrid bundles \\
BLS12-381 & checkpoints, escrow, arbitration, credit attestations & quorum
  certificates and threshold aggregation \\
ML-DSA-65 & checkpoint sidecar, post-quantum mode & post-quantum signatures
  without pairing aggregation \\
Argon2id & wallet file & passphrase-based key derivation \\
ChaCha20-Poly1305 & wallet file, node state container & authenticated
  encryption \\
\bottomrule
\end{tabular}
\caption{Primitive inventory. The BLS implementation is a reference pairing
library with published audit history \cite{supranational-blst}; the aggregate
signature scheme follows the standard construction and its known rogue-key
attacks \cite{boneh2003aggregate,boneh2018multisignatures}.}
\label{tab:crypto}
\end{table}

Because BLS aggregation has no pairing-based analogue in the lattice setting,
the post-quantum signature travels in a sidecar: a bounded list of per-signer
lattice signatures accompanying the checkpoint. A signer bitmap that holds a
registered lattice key but presents no sidecar entry is refused --- a sidecar
stripped in transit lands as ``missing'', never as ``valid without''.

\subsection{Consensus-level security}

\paragraph{Integer-only arithmetic.} Every consensus computation is integral.
Difficulty, supply, fees and the proof-of-work forward pass all refuse
overflow rather than wrapping it, and the release profile enables overflow
checks in the compiler.

\paragraph{Ordering of checks.} A submitted block is checked for parent
existence and for agreement with the difficulty its own branch demands
\emph{before} the proof of work is verified. This ordering is not cosmetic: a
header that declares a trivial difficulty has a trivially satisfiable target
but still costs the verifier a full proof check, so checking the difficulty
after the proof would let anyone burn the node's CPU for free, repeatedly. The
fix was applied across the consensus, peer-to-peer, API, command-line and
Model Context Protocol surfaces in a single audit pass.

\paragraph{Reorganisation bounds.} An attempted reorganisation that would cross
the binding finalised point is refused, and the node state is restored.

\paragraph{Proof-of-work diversity.} The alternative memory-hard puzzle is
specified, implemented, separately reviewed and left unregistered, so it
cannot affect consensus until it is deliberately enabled.

\subsection{Network and API security}

The peer-to-peer layer bounds every quantity an attacker can inflate: frame
payload size, inventory items per message, transaction request batch size,
in-flight and pending relay items, remembered transactions, per-peer per-minute
transaction and control budgets, total and inbound peer counts, in-flight
handshakes, and the number of arbitration tallies remembered.

The HTTP API enforces, on every bind: a request-body size limit; a cap on
concurrent requests in flight; a per-request deadline; a separate deadline for
the client to finish sending request headers, which is the slowloris bound; a
connection cap enforced in a manual accept loop so that an over-cap socket is
dropped before a single byte is read; a host-header check that refuses a
missing, unknown or bracket-malformed host on every bind including loopback,
which is a defence against DNS rebinding; and, on any non-loopback bind, a
mandatory bearer token compared in constant time over hashes so that neither
the value nor its length leaks through timing. Non-loopback binds are refused
outright if no token is configured. Error responses are uniformly JSON with a
stable machine-readable code and a human-readable detail string, and the detail
never contains signing material. Rate limiting is a token bucket per remote
address, with IPv6 addresses grouped by prefix.

\subsection{Replay protection and fail-closed behaviour}

\begin{table}[t]
\centering
\begin{tabular}{lll}
\toprule
\textbf{Surface} & \textbf{Guard} & \textbf{Behaviour on replay} \\
\midrule
Payment & exact-payment tracker, capped & refuse (never evict) \\
Receipt & receipt book, contradiction detection & refuse by name \\
Credit attestation & consumed-attestation set, persisted & refuse as reuse \\
Credit reservation & reservation log, idempotent & no second reservation \\
Policy authorisation & strictly increasing nonce & refuse \\
Settlement link & one payment, one transaction & refuse as already settled \\
Escrow-to-exit binding & one escrow, one exit & refuse as already settled \\
Exit identity & $(\text{owner},\text{nonce})$ consumed set & release exactly
  once, ever \\
State file & SHA-256 digest, fail-closed load & refuse to start \\
\bottomrule
\end{tabular}
\caption{Replay guards. The general policy is refusal rather than eviction:
dropping the oldest entry would silently re-open an old replay window.}
\label{tab:replay}
\end{table}

\subsection{The lying-root rule}
\label{sec:lyingroot}

One security property deserves a dedicated subsection because it is unusual and
because it demonstrates the design's threat model.

The Layer-2 state root committed by a block is checked for \emph{presence},
not for correctness. A block that commits a root inconsistent with the
Layer-2 state the proposer actually computed is accepted at Layer 1. The
consequence is bounded and one-way: exits whose inclusion proofs do not verify
against the committed root simply never become settleable, so they stall.
They do not pay the wrong party, they do not pay twice, and they do not corrupt
the Layer-1 ledger. Splitting proof-checking from chain-checking is what makes
a forged proof cryptographically impossible even when the block producer
commits a lying root.

This is a deliberate trade. Requiring every node to re-execute the entire
Layer-2 state machine to validate a block would make Layer 2 a consensus
layer, with all the latency and reorganisation costs that implies. Instead, the
cost of a dishonest producer is paid by the honest users of the exit queue
whose exits stall --- and by the producer, whose own exits stall too. A test
suite mines blocks past the finality depth to demonstrate that the stall is
permanent rather than merely immature.

\subsection{Security engineering as a process}
\label{sec:security-process}

Security in this project is treated as a measurable process rather than a
claim.

\paragraph{Four parallel audits.} A dedicated audit pass covered the API, the
wallet and command-line surfaces, the peer-to-peer layer, and the Layer-2 and
agent surfaces in parallel, producing a set of fixes applied together. The
findings were, in outline: the consensus check-ordering problem described
above; an unbounded peer-event queue replaced with a bounded one that pushes
backpressure onto the sender instead of buffering; the in-flight handshake cap;
per-peer rate limits on control messages, which are pairing-checked on arrival
and would otherwise burn CPU in a flood; a cap on remembered arbitration
tallies; a capacity check moved ahead of deduplication scans in the inventory
handlers; length-hiding token comparison; IPv6 rate-limit buckets keyed by
prefix; refusal of empty-token public binds; saturating arithmetic in the
transaction builder; an off-by-one in the spendable-balance calculation; path
traversal through a validator key label; a queue drain that only happened while
mining; refusal to overwrite an existing wallet file; and errors during
shutdown flush that were previously discarded.

\paragraph{Red-team suites.} Four adversarial test files target the consensus
and ledger rules directly --- twenty-one tests exercising attempted
misbehaviour against the chain, the credit path, the delegation path and a
second round of edge cases.

\paragraph{Fuzzing.} Seventeen cargo-fuzz targets cover the attacker-controlled
boundaries: transaction decoding, block decoding, wire-message decoding,
mempool admission, storage recovery, checkpoint validation and codec, metadata
parsing, Layer-2 payload decoding, Layer-2 state application, the dispute
arbitration engine, verified work certificate decoding, the machine-readable
task readers, the deliverable codec, the worker decision flows, and two
harnesses written for the authority and credit layers: one that drives grants,
delegations, revocations and spends, and one that drives loans, collateral,
repayments and defaults. The last two assert the properties those layers
depend on --- value conservation, that a spend never exceeds its grant, that
revocation is absolute, that authority narrows, that debt only goes one way,
that a default is not early, and that a replay is refused. A reported campaign
of 9\,352\,161
executions produced no crashes and no sanitizer findings
\cite{jounit-fuzz}; every registered target is now run in the continuous
integration smoke loop, and every fuzz target has a deterministic counterpart in
the ordinary test suite, so the properties the fuzzer checks are also checked on
every CI run.

\paragraph{An independent implementation of the authority and credit rules.}
The rules of Section~\ref{sec:authority} are implemented a second time, in a
different language, written from the specification rather than from the
protocol's source. It derives the addresses, the agent identities, the grant
hashes and the position identifiers itself, and emits a vector file of
scenarios; the protocol's test suite replays those scenarios and asserts the
same verdict, the same reason and the same final state. This is the same
discipline the proof-of-work reference uses, applied to the part of the
protocol where a rule --- rather than a digest --- is the thing that has to
agree.

\paragraph{Attack simulation.} Independent Python simulations reproduce the
difficulty-adjustment behaviour under adversarial timing, the tail-emission
schedule, and the puzzle digests, and a cross-check script verifies the
consensus vectors against the Rust implementation.

\paragraph{Disclosure posture.} The readiness record documents a ninety-day
responsible-disclosure window, a severity model in which a chain-halting defect
is critical and a key compromise is high, and a contact address --- while
noting honestly that the corresponding public key is not yet published, so the
channel is not yet operational. We consider stating that plainly to be more
useful than presenting the policy as if it were complete.

\subsection{Known gaps}
\label{sec:gaps}

The following are known and are stated rather than concealed.

\begin{itemize}
\item No independent third-party security audit has been completed. The
readiness record says so, and the recommendation stands: a third-party review
should precede any deployment carrying real value \cite{jounit-readiness}.
\item The disclosure contact exists but its public key is unpublished, so the
channel is not yet operational.
\item The specification and the code are not in perfect agreement everywhere;
several drifted sections were corrected, but the specification remains a
living document and the code is authoritative where they disagree in
practice.
\item The peer-to-peer layer has no discovery protocol and no peer
authentication, and the Layer-2 dispute window admits a payment of an
unchallenged result that happens to be incorrect.
\item The network has not been subjected to load, eclipse or economic testing
at mainnet scale.
\item There is no validator yield yet. The decision is recorded --- the reward
comes from the Layer-2 fee pool, so the Layer-1 coinbase is untouched --- but
it is not implemented, because paying it means deciding which checkpoint a
committee's work counts for, and the checkpoint a committee signs is not part
of a block.
\item There is no exchange: no second asset, no order book, no price oracle.
The intent is that JNT's economic usage comes from JNT network usage.
\end{itemize}


\section{Use by humans, machines, robots and agents}
\label{sec:use}

\subsection{Identities and roles}

The protocol defines three roles. A \emph{human} is an ordinary wallet. A
\emph{machine} is a device spending under a policy written by its owner. A
\emph{service} is an identity that receives payments. The agent identifier is
derived from the network identifier, the role and the owner, so the same device
has different identifiers on different chains --- a deliberate property, since
an identity that is portable across chains is an identity that can be
accumulated against on one chain and redeemed on another.

\subsection{What an agent may do}

An agent can hold value, receive value, spend value under a written grant,
publish and revoke grants of its own, lend to another agent and be lent to,
post and claim tasks, execute work, submit attested results, challenge
results, and settle or be settled. What an agent cannot do is exceed the grant
it was given, replay an authorisation, act without an owner signature
somewhere in its ancestry, borrow without a lender's signature, or cause the
protocol to create value. There is no
second currency: the owner field is an ordinary address, so machine and human
balances live in the same output set under the same issuance schedule.

\subsection{Agent tooling}

\paragraph{Worker SDK.} The worker crate is a client, not a consensus
component. It discovers open tasks, claims them, executes them, submits
results, raises disputes and requests exits, and its decision functions are
deterministic --- no clock, no randomness --- so the same inputs always produce
the same plan. A dedicated security suite asserts that an actor starting with
zero privilege cannot, through any path in the SDK, become an authorised
worker, claimant, verifier, settler or payment recipient. A second module covers
the authority and credit flows: publish a grant, spend under it, revoke it,
open a loan, add collateral, repay, and settle a default --- each one planned
locally and then broadcast, like every other path in the SDK.

\paragraph{Model Context Protocol surface.} The MCP crate is a keyless
dispatcher over agent identity, payments, receipts and reputation, and over the
authority and credit layer: an operation reads a published grant, a spend under
a grant, or a loan, and returns the terms and the hash to sign. It is
explicitly \emph{not} a trust authority: it verifies using the same public
functions every other verifier uses, and it cannot mint an authorisation, a
payment or a receipt. It is explicitly \emph{not} a key holder: no operation
accepts, returns or stores a private key; where a signature is needed it
returns the bytes and the signing hash, and the client signs. It holds an
authorisation book so that anti-replay counters survive across calls, because
a fresh book per call would silently reset every counter. Its transport is a
small dependency-free HTTP client.

\paragraph{The command line.} The node's own command line carries the same
flows, because an operator should not have to write a program to grant an
agent a budget: a deposit address and a Layer-2 summary, a grant publication,
a spend, a revocation, and the credit operations --- each of which checks the
transaction against the state the chain will replay and refuses with the
chain's own reason, rather than mining a block the replay would skip.

\subsection{Concrete scenarios}

\paragraph{A research pipeline buys a computation.} A submitter posts a task
with an escrow and a compute specification. A worker claims it, runs it, and
submits a result hash with a signed attestation. If nobody challenges the
result within the dispute window, the task settles correct and the worker is
paid from the escrow; a verified work certificate is issued, and a relying
party may insist on the quorum-backed kind if it needs independent proof.

\paragraph{A robot fleet operates under a budget.} An operator publishes a grant
bounding an agent per transaction, bounded per window, restricted to named
recipients and named services, with an increasing nonce. The fleet spends
within that envelope, and every node enforces it, not only the one that issued
it. The operator revokes the grant, and every descendant grant dies with it.

\paragraph{A machine buys a service from another machine.} A service identity
publishes a receipt; the payer records the payment; the payment is settled
only when bound to a Layer-1 transaction that actually pays the recipient. The
receipt and the payment cannot be confused with one another, and a
contradicting receipt for the same payment is refused.

\paragraph{A human buys the same service.} The same paths apply. There is no
separate human lane, no privileged account type, and no way for the protocol
to distinguish a paying human from a paying robot other than by the identity
they present.


\section{Openness, governance and control of information}
\label{sec:governance}

\subsection{Licence and ownership}

The code is released under a standard permissive Rust dual licence,
\texttt{MIT OR Apache-2.0}, with both texts at the repository root
\cite{jounit-repo}. Anyone may run it, study it, modify it, audit it,
redistribute it, run a public node or a fork, use it as a dependency, write
about it, and report vulnerabilities against it. Commercial use is permitted by
the licence; the trademark document asks for a short written permission for
\emph>commercial use of the marks themselves\}, which is a different question
from commercial use of the code.

The project states that it is community-owned in the Bitcoin sense, and that no
single person holds it. One inconsistency should be disclosed: the MIT licence
text names an individual copyright holder, while the project's own
documentation asserts community ownership. Reconciling those two statements is
an open administrative task.

\subsection{Trademark separation}

The name, the ticker and the logos are community marks. Permitted without
asking: running the software on any network, forking it, using the marks
factually to refer to the project (``built on Jounit'', in a compatibility list,
in an article). Not permitted without asking: naming a product, company, chain
or service ``Jounit'' or ``JNT''; using the logos in one's own branding;
publishing a fork under the name or implying official affiliation; selling
something under the name; registering the marks as one's own.

The reasoning is worth stating because it is a general principle about
permissively licensed infrastructure: a permissive code licence cannot protect
a brand, and pretending otherwise fails. The code stays open; the name stays
with the community; a fork takes a new name because two chains both displaying
\texttt{JNT} is exactly the confusion that harms users.

\subsection{Governance after launch}

The project's stated principle is no post-launch governance: no foundation
controls the protocol, no upgrade path is reserved for an insider, and the
rules are the rules. Difficulty and puzzle parameters are already designed to
be adjustable, but the supply ladder, the halving schedule and the base-fee
split are stated to be immutable, so monetary policy cannot be changed by a
software update. This is the Bitcoin-style commitment to a rule rather than to
a committee.

\subsection{Who holds and controls the information}

\begin{table}[t]
\centering
\begin{tabular}{@{}p{0.24\textwidth}p{0.29\textwidth}p{0.38\textwidth}@{}}
\toprule
\textbf{Information} & \textbf{Where it lives} & \textbf{Who controls it} \\
\midrule
Ledger, transactions, task
  specifications & On chain, replicated & Nobody individually; rule-governed \\
Task results and attestations & On chain & Nobody individually \\
Exit records & Layer-2 state, committed by
  Layer-1 root & Nobody individually \\
Published grants, their counters & Layer-2 state, committed by
  the Layer-1 root & The owner, within the limits the grant states \\
Credit positions and their collateral & Layer-2 state, committed by
  the Layer-1 root & The lender and the borrower, by agreement \\
Wallet private keys & The holder's file, encrypted & The holder alone \\
Validator signing keys & The validator's encrypted file & The validator
  alone \\
Credit issuer key & The operator's configuration & The operator \\
Receipts, payment links, service-layer
  credit and escrow books & Node-local state file & The node operator, per node \\
Reputation & Derived view of chain events & Nobody; derived, not owned \\
Delegation lineage records & Node-local, signed by the owner & The owner,
  subject to revocation \\
\bottomrule
\end{tabular}
\caption{Information-control matrix. The distinction that matters is between
information that is \emph{committed to the chain} and therefore replicated and
immutable, and information that is \emph{node-local} and therefore belongs to
one operator.}
\label{tab:info}
\end{table}

Two properties follow from this table. First, no single party holds the
information needed to reconstruct the global state of the system: the ledger is
replicated, and node-local books are explicitly \emph{not} an alternative
source of truth for ownership or for spendable balance. A corrupt local book
can cost an operator their own bookkeeping; it cannot mint a coin or move one.
Second, the system is asymmetric on purpose: the operator's own books are
private to them, and the shared financial record is public to all and owned by
none. This is the same asymmetry that makes an open ledger useful for audit
and a private wallet file useful for safety.

The node-local boundary is enforced by construction rather than by policy: the
receipt, payment, service-layer credit and escrow books are not reachable from
the consensus path, and the documented rule is that these books must never be
treated as an alternative source of truth for global ownership or for a
spendable balance. Coin movement happens only through Layer-1 outputs and
Layer-2 account balances. The spending grants and the credit positions are on
the other side of that line: they are consensus state, every node holds them,
and they are covered by the state root.


\section{Using Jounit}
\label{sec:usage}

\subsection{Building and running}

The implementation requires Rust 1.75 or later (developed and tested against a
much newer release) and Python 3.11 or later for the simulation suite.

\begin{verbatim}
git clone https://github.com/swamy18/Jounit.git
cd Jounit
cargo build --release
\end{verbatim}

A chain directory, a miner wallet and the genesis block are created with
\texttt{init}; blocks are produced with \texttt{mine}; the node reports its tip,
difficulty, supply and balances with \texttt{status}; and \texttt{verify}
replays the whole chain from genesis and diffs the result against stored state.
\texttt{serve} runs the peer-to-peer node, mines by default so that matured
exits settle, and optionally exposes the HTTP API.

\begin{verbatim}
export JNT_WALLET_PASSPHRASE=change-me
jnt-node init  --network testnet --data-dir ./chain
jnt-node mine  --network testnet --data-dir ./chain --blocks 25
jnt-node verify --network testnet --data-dir ./chain
jnt-node serve --network testnet --data-dir ./chain --listen 127.0.0.1:9833
\end{verbatim}

For the HTTP API, a loopback bind needs no credentials; a non-loopback bind
requires a bearer token, supplied through the environment rather than a flag,
because flags are visible in the process table.

\begin{verbatim}
export JNT_API_TOKEN="..."
jnt-node serve --network testnet --data-dir ./chain --api-listen 0.0.0.0:8080
\end{verbatim}

\subsection{Command surface}

\begin{table}[t]
\centering
\begin{tabular}{ll}
\toprule
\textbf{Command} & \textbf{Function} \\
\midrule
\texttt{init} & Create a chain directory, a miner wallet and the genesis block \\
\texttt{mine} & Mine blocks and extend the chain \\
\texttt{send} & Spend from a wallet address, locally or over the network \\
\texttt{status} & Chain tip, difficulty, supply, balance, Layer-2 pipeline \\
\texttt{verify} & Replay the chain from genesis and diff it against storage \\
\texttt{show-block} & Print one block \\
\texttt{balance} & List unspent outputs of an address \\
\texttt{wallet} & Create or list wallets \\
\texttt{policy} & Publish, revoke or list spending grants \\
\texttt{agent-spend} & Spend under a published grant \\
\texttt{credit} & Open, repay, default or list credit positions \\
\texttt{l2} & Layer-2 summary, and the deposit address \\
\texttt{validator} & Add or list checkpoint-signing keys \\
\texttt{export} & Write the canonical chain to JSON \\
\texttt{serve} & Run the node: relay, sync, API, exit crank, signing \\
\bottomrule
\end{tabular}
\caption{The node's command surface. Every command accepts \texttt{--help}, and
every command accepts a network profile and a data directory.}
\label{tab:cli}
\end{table}

\subsection{Authority and credit, from the terminal}

A grant is the owner's signed authority for a named spender; it is published on
chain, so every node enforces it. A loan is the same idea for value.

\begin{verbatim}
# The deposit address: an ordinary Layer-1 payment to it credits a
# Layer-2 balance, which is what the grants and the loans spend from.
jnt-node l2 --deposit-address
jnt-node send --to <deposit-address> --amount 5

# Publish a grant: this key may spend up to 0.01 JNT per payment and
# 0.05 JNT per window, to this recipient only.
jnt-node policy --spender <spender> --per-tx 1000000 \
                --cumulative 5000000 --recipients <recipient>

# The agent spends under it.
jnt-node agent-spend --policy <policy-hash> --to <recipient> \
             --amount 0.002 --nonce 1

# The owner kills it; every descendant of the grant dies with it.
jnt-node policy --revoke --agent-id <agent-id>
jnt-node policy --list
\end{verbatim}

\begin{verbatim}
# A loan: the lender signs, the principal moves, the collateral answers
# for the debt.
jnt-node credit --to <borrower> --amount 1 --collateral 100000000
jnt-node credit --repay --position <position-hash> --amount 1
jnt-node credit --list
\end{verbatim}

Every one of these commands checks the transaction against the state the chain
will replay and refuses with the chain's own reason, rather than mining a block
the replay would skip.

\subsection{Profiles}

The three profiles differ only in numbers, and the \texttt{mini} profile is
the same code path as production.

\begin{table}[t]
\centering
\begin{tabular}{llll}
\toprule
\textbf{Parameter} & \textbf{mainnet} & \textbf{testnet} & \textbf{mini} \\
\midrule
Block time & 600\,000 ms & 1\,000 ms & 1\,000 ms \\
Halving period & 262\,800 & 2\,000 & 2\,000 \\
Start reward & 76 JNT & 76 JNT & 76 JNT \\
Finality depth & 100 & 20 & 20 \\
Coinbase maturity & 100 & 10 & 10 \\
Genesis validators & none & 4 & 4 \\
Binding finality & height 100 & height 20 & off \\
Exits & height 1 & height 1 & height 1 \\
Post-quantum requirement & off & off & off \\
\bottomrule
\end{tabular}
\caption{Network profiles. The mainnet profile deliberately ships with no
genesis committee, so that every checkpoint fails closed until the
stake-weighted committee derives from the bond ledger --- the protocol never
verifies against nothing.}
\label{tab:profiles}
\end{table}


\section{Limitations and future work}
\label{sec:future}

\subsection{Limitations}

The limitations of the present system are structural rather than incidental,
and it is better to state them than to let a reader infer them.

\paragraph{The consensus layer is conventional.} The fork choice is cumulative
work on a linear chain. This buys predictability and a well-understood security
model at the cost of throughput and of the latency properties that a DAG or a
committee-based design would offer \cite{eyal2014securitybitcoin,trivedi2017sok}.
Jounit's ambition is to be a settlement layer, and the design reflects that
choice rather than transcending it.

\paragraph{Layer 2 inherits Layer 1's ordering and none of its latency.} The
replay model means Layer-2 state is always as fresh as the last Layer-1 block,
and a Layer-2 payload that fails to apply is skipped rather than reverting the
block. This is a deliberate trade of strictness for liveness.

\paragraph{The optimistic settlement model has a known cost.} An unchallenged
incorrect result is paid. The protocol narrows the consequence --- a
quorum-backed certificate is distinguishable from an unchallenged one --- but
does not remove it.

\paragraph{Peers are unauthenticated.} The current transport has no node
identity and no discovery protocol. A denial-of-service or eclipse attack
against the peer-to-peer layer is mitigated by bounds and bans but not
prevented by cryptography.

\paragraph{The Layer-2 dispute machinery depends on a committee.} Committee
membership is derived from bonded stake and capped at twenty-one members, so
the adjudicator is small. That is a deliberate consequence of quorum
certificates, and it concentrates the value of a successful compromise.

\paragraph{No field deployment.} The chain has not been run at mainnet scale,
under load, or against an adversarial peer-to-peer network. The verification
record's benchmarks are single-machine measurements on modest hardware, and the
fuzzing campaign is bounded rather than exhaustive.

\subsection{Future work}

In rough dependency order: an independent third-party security review and a
published disclosure key; a peer-to-peer layer with node identity, peer
authentication and a discovery mechanism; a larger committee with staggered
rotation; a richer proof-of-work registry in which a memory-hard puzzle can be
enabled by consensus rule rather than by configuration; a formal verification
effort on the consensus transition function and the proof-of-work forward
pass; and a long-horizon economic study of the issuance schedule under
realistic fee activity, which would test the claim that net supply can fall.


\section{Conclusion}
\label{sec:conclusion}

Jounit is an attempt to build a settlement primitive that general-purpose
cryptocurrencies do not provide: a way for a machine to be paid for work that a
third party can verify without re-running it, and for a principal to delegate
spending authority to such a machine without losing the ability to bound it.

The design is not exotic. It uses a linear UTXO chain with cumulative-work
fork choice, a quorum certificate over a stake-weighted committee, a
deterministic replay model for the second layer, and inclusion proofs to move
value back across the layer boundary. What is unusual is the set of constraints
applied throughout: no floating point in consensus, overflow refused rather
than wrapped, a work-share band that prevents the interesting part of the
proof of work from becoming the weak part, an escrow that names its own
arbitrator and refuses oracle conditions it cannot evaluate, a reputation
system that is structurally excluded from consensus, a receipt system that
distinguishes ``nobody objected'' from ``a quorum verified this'', a
presence-only rule for Layer-2 roots that converts a dishonest producer's
power to steal into a power to stall, and a wallet file that is encrypted,
atomically written and zeroised.

The newest part of the design is the least exotic and, we would argue, the most
consequential: a machine's authority to spend, and a machine's ability to
borrow, are stated as \emph{rules the chain checks} rather than as records one
node keeps. The grant binds an owner to a named spender and a named agent; the
loan binds a lender to a borrower, a principal, a collateral, a rate and a
deadline. Neither is novel as an idea. What is new is that they were
node-local, and that closing that gap was mostly a matter of writing the rules
down twice --- once in a specification, once in a second implementation --- and
letting the two disagree until they did not.

The honest summary of the security posture is that the engineering is
disciplined and the evidence is unusually well documented for a pre-launch
project --- a measured fuzzing campaign, red-team suites, independent
re-implementations, a spec-versus-code reconciliation pass, and a
disclosure of what was previously overstated --- while an independent
third-party review has not yet been performed. That is stated in the paper and
in the project's own records because a protocol's credibility ultimately rests
on the accuracy of its self-assessment.

The vision is a chain that does useful work with the energy it consumes, and
that treats an autonomous agent as a first-class economic participant rather
than as an account to be impersonated. Whether that vision survives contact
with adversarial use is an open question, and the correct next step is the one
that has not yet been taken.


\section*{Acknowledgments}
The author thanks the contributors to the open-source cryptocurrency and
agent-runtime ecosystems whose prior work this protocol builds on, and the
reviewers of the specification whose disagreements improved it. The
specification, the decision log and the verification record are the primary
artefacts of this work and are openly available
\cite{jounit-spec,jounit-decisions,jounit-status,jounit-readiness,jounit-monetary,
jounit-fuzz}.


\bibliographystyle{unsrt}
\bibliography{references}


\appendix

\section{Key protocol constants}
\label{app:constants}

\begin{table}[t]
\centering
\begin{tabular}{llr}
\toprule
\textbf{Quantity} & \textbf{Where} & \textbf{Value} \\
\midrule
Block header size & block & 177 bytes \\
Transaction output size & block & 28 bytes \\
Address size & hash & 20 bytes \\
Agent identity size & agent & 20 bytes \\
Grant hash size & agent & 32 bytes \\
Gas per transaction & transaction & 16\,000 \\
Gas per input & transaction & 1\,000 \\
Gas per output & transaction & 250 \\
Maximum transactions per block & parameters & 5\,000 \\
Maximum block size & parameters & 4\,000\,000 bytes \\
Difficulty window & parameters & 60 blocks \\
Maximum difficulty change & parameters & 10\,\% \\
Median look-back & difficulty & 11 blocks \\
Maximum attempts per block & chain & 200\,000\,000 \\
Quorum threshold & validators & $\lceil 2n/3 \rceil$, basis points 6\,667 \\
Maximum committee size & validators & 21 \\
Epoch length & validators & 100 blocks \\
Unbonding period & bond & 1\,008 blocks \\
Slashing vesting period & slashing & 52\,560 blocks \\
Escrow expiry & escrow & 1\,008 blocks \\
Layer-2 deposit minimum & parameters & 1\,000\,000 subunits \\
Layer-2 flat fee minimum & parameters & 1\,000 subunits \\
Exit fee minimum & exit & 1\,000 subunits \\
Whistleblower bounty & task contract & 25\,\% \\
Layer-2 transaction codes & layer 2 & \texttt{0x10}--\texttt{0x1C} \\
Grant / credit leaf tags & state root & \texttt{0x05} / \texttt{0x06} \\
Delegation depth cap & delegation & 8 \\
Grant publication fee & parameters & 2\,000 \\
Grant revocation fee & parameters & 2\,000 \\
Agent spend fee & parameters & 2\,000 \\
Credit open / collateral / repay / default fee & parameters & 2\,000 each \\
Maximum credit rate per window & parameters & 1\,000 bps \\
Maximum frame payload & network & 4\,100\,000 bytes \\
Maximum frame payload & network & 4\,100\,000 bytes \\
Maximum peers (total / inbound) & network & 125 / 32 \\
Maximum API connections / concurrency & API & 128 / 64 \\
API request timeout / header timeout & API & 10 s / 30 s \\
Token bucket rate / capacity & API & 60 per second / 120 \\
Maximum state container size & store & 64 MiB \\
\bottomrule
\end{tabular}
\caption{Selected constants. Every value is a named constant in the source
rather than a literal at a use site, and every profile derives from the same
parameters module \cite{jounit-spec}.}
\label{tab:constants}
\end{table}

\section{Security invariants}
\label{app:invariants}

The following invariants are asserted by the test suite, the fuzz targets, or
both.

\begin{enumerate}
\item A block that does not satisfy the proof of work, or whose declared
  difficulty disagrees with its branch, is refused --- and the difficulty
  check precedes the proof check.
\item No transaction may create value: the sum of inputs equals the sum of
  outputs plus the fee.
\item No block may mint beyond the issuance integral or the rate anchor.
\item The base fee is split exactly in half between burn and reserve, with the
  odd unit assigned to the burn.
\item A reorganisation that would cross the binding finalised point is refused.
\item A checkpoint is accepted only if it satisfies every clause of the
  seven-clause checklist, including quorum, signature and sidecar.
\item A layer-2 state transition preserves the conservation invariant.
\item An exit may be released at most once per (owner, nonce) pair, and only
  to the destination recorded in the proven exit record.
\item A settlement link binds one payment to at most one transaction, and the
  transaction must pay the recipient.
\item A credit attestation may mint at most once, ever, including across a
  restart.
\item A spending authorisation cannot be replayed: the nonce must strictly
  increase.
\item A destroyed state file is a refusal to start, never an empty book.
\item A lying layer-2 root is accepted at layer 1 but arms no exit.
\item The API never signs, never holds a private key, and never reaches the
  consensus submission path.
\end{enumerate}

\end{document}
